% f06_compiler variant c -- the obligations as a dependency graph: what each
% one rests on, and which pass can establish it.
% Data: src/sprs_core.py -- CBTree (S1, l.395-497), _validate_assignment
%       (l.1424), DMEMAllocator.allocate (l.5143-5177), the operand lookup in
%       build_schedule/generate_instructions (main.tex quotes the line),
%       Topology.generate_routing_tables (l.681-700) and its use in
%       emit_sv_testbench (l.5602-5610). Obligation statements and the
%       hypothesis lists of Claim 1 and Property 2: main.tex, App. A.
%       No counts appear in this variant; it is a structure, not a measurement.
% Strategy: the ORDER of the argument rather than the order of the passes.
% It is the only variant that shows what each obligation rests on -- and that
% nothing rests on the mapping, which is Proposition (orthogonality) drawn
% rather than asserted.
% Colour: the argument runs left to right and changes colour as it changes
% kind -- spBlue for the inputs and for the passes that ran, spGreen for the
% obligations those passes discharge (the same green as the "establishes"
% column of f04 variant b), spInk for the two conclusions. The one edge that
% carries nothing is spMute and dashed: mu reaches the emitter only through NOP
% padding, so the figure's claim of orthogonality is visible as the single grey
% line in a blue graph. The dagger on C-route is spAmber, the paper's caution
% mark. Dash pattern, rule weight and column position carry all of it with the
% colour stripped.
% Target: IEEE double column, 7.16in = 18.2cm. Drawn inside 17.0cm, which is
%         also the standalone wrapper's varwidth, so nothing is clipped.
\begin{tikzpicture}[x=1cm, y=1cm,
  inp/.style={draw=spBlue, line width=0.4pt, rounded corners=2pt, fill=white,
             font=\tiny, align=center, inner sep=2.5pt, text width=20mm,
             text=spInk},
  pa/.style={draw=spBlue, line width=0.55pt, rounded corners=2pt, fill=spBlueF,
             font=\tiny, align=center, inner sep=2.5pt, text width=30mm,
             text=spInk},
  ob/.style={draw=spGreen, line width=0.7pt, rounded corners=2pt, fill=spGreenF,
             font=\tiny, align=center, inner sep=2.5pt, text width=32mm,
             text=spInk},
  cl/.style={draw=spInk, line width=0.7pt, rounded corners=2pt, fill=white,
             font=\tiny, align=center, inner sep=2.5pt, text width=26mm,
             text=spInk},
  e/.style={-{Latex[length=1.5mm,width=1.1mm]}, draw=spBlue, line width=0.4pt},
  ed/.style={e, dash pattern=on 1.4pt off 1.2pt, draw=spMute},
  hd/.style={font=\scriptsize\bfseries, anchor=south, text=spInk},
]
\node[hd] at (1.30,0.85) {inputs};
\node[hd, text=spBlue] at (5.60,0.85) {the pass that establishes it};
\node[hd, text=spGreen] at (10.60,0.85) {obligation};
\node[hd] at (15.20,0.85) {used by};

% ---- inputs -----------------------------------------------------------
\node[inp] (N)   at (1.30,-0.45) {$N$: $\FBT(N)$ is a tree --- every node has exactly one consumer};
\node[inp] (A)   at (1.30,-2.20) {$\alpha$ from S2, a valid partition of the nodes};
\node[inp] (M)   at (1.30,-3.85) {$\mu$ from S3};
\node[inp] (T)   at (1.30,-5.30) {$(G,\tau)$: the adjacency table};

% ---- passes -----------------------------------------------------------
\node[pa] (S1)  at (5.60, 0.05) {S1 fixes the leaf and operand binding};
\node[pa] (AL)  at (5.60,-1.85) {\sv{DMEMAllocator.allocate}: SU post-order, free list, then a never-freed receive-buffer pass};
\node[pa] (EM)  at (5.60,-3.85) {the emitter's operand lookup, \sv{addr\_map[gpu][c]}};
\node[pa] (RT)  at (5.60,-5.30) {\sv{generate\_routing\_tables}};

% ---- obligations ------------------------------------------------------
\node[ob] (Cleaf)  at (10.60, 0.05) {\textsf{C-leaf} the leaf a \sv{GENERATE} writes is the one $\Rcan$ names};
\node[ob] (Csw)    at (10.60,-1.55) {\textsf{C-sw} one writer per receive buffer, and it is remote};
\node[ob] (Cmem)   at (10.60,-2.85) {\textsf{C-mem} within \sv{DMEM\_DEPTH} $=\DmemDepth$ and the router-id sentinel};
\node[ob] (Cbind)  at (10.60,-4.15) {\textsf{C-bind} a \sv{COMPUTE}'s operands are its node's two children};
\node[ob] (Croute) at (10.60,-5.45) {\textcolor{spAmber!80!spInk}{$\dagger$}\,\textsf{C-route} destination-correct, loop-free paths};

% ---- consumers --------------------------------------------------------
\node[cl] (C1) at (15.20,-1.55) {\textbf{Claim 1} bitwise invariance --- partial correctness};
\node[cl] (P2) at (15.20,-4.60) {\textbf{Property 2} presence-gated fail-stop --- liveness};

% ---- edges ------------------------------------------------------------
\draw[e] (N) -- (S1);
\draw[e] (N.east) -- ++(0.35,0) |- (AL.west);
\draw[e] (A) -- (AL);
\draw[e] (A.east) -- ++(0.55,0) |- (EM.west);
\draw[e] (T) -- (RT);
\draw[ed] (M.east) -- ++(0.75,0) |- (EM.south);
\draw[e] (S1) -- (Cleaf);
\draw[e] (S1.east) -- ++(0.30,0) |- (Cbind.west);
\draw[e] (AL) -- (Csw);
\draw[e] (AL.east) -- ++(0.30,0) |- (Cmem.west);
\draw[e] (EM) -- (Cbind);
\draw[e] (RT) -- (Croute);
\draw[e] (Cleaf.east) -- ++(0.35,0) |- ([yshift=2.6mm]C1.west);
\draw[e] (Csw) -- (C1);
\draw[e] (Cmem.east) -- ++(0.35,0) |- ([yshift=-2.6mm]C1.west);
\draw[e] (Cbind.east) -- ++(0.60,0) |- (C1.south west);
\draw[e] (Cbind.east) -- ++(0.60,0) |- (P2.north west);
\draw[e] (Croute.east) -- ++(0.35,0) |- (P2.west);

% ---- the two things this graph shows ---------------------------------
\node[font=\tiny, anchor=north west, align=left, text width=78mm, text=spInk]
  at (0.10,-6.05)
  {\textcolor{spMute}{\emph{Nothing depends on $\mu$.}} Its only edge is dashed and reaches the emitter
   through the cost model's NOP padding; no obligation, and neither Claim~1 nor
   Property~2, has it as a hypothesis.};
\node[font=\tiny, anchor=north west, align=left, text width=82mm, text=spInk]
  at (8.60,-6.05)
  {\textcolor{spGreen}{\emph{One walk carries two obligations.}} \textsf{C-sw} is \emph{constructed} by the
   allocator's post-order rather than assumed, and \textsf{C-mem} falls out of the same
   walk as an $O(N)$ read-off --- which is why S5 never had to run.};
\end{tikzpicture}
